%% Beginning of file 'sample7.tex'
%%
%% Version 7. Created January 2025.  
%%
%% AASTeX v7 calls the following external packages:
%% times, hyperref, ifthen, hyphens, longtable, xcolor, 
%% bookmarks, array, rotating, ulem, and lineno 
%%
%% RevTeX is no longer used in AASTeX v7.
%%
\documentclass[twocolumn]{aastex7}
%%
%% This initial command takes arguments that can be used to easily modify 
%% the output of the compiled manuscript. Any combination of arguments can be 
%% invoked like this:
%%
%% \documentclass[argument1,argument2,argument3,...]{aastex7}
%%
%% Six of the arguments are typestting options. They are:
%%
%%  twocolumn   : two text columns, 10 point font, single spaced article.
%%                This is the most compact and represent the final published
%%                derived PDF copy of the accepted manuscript from the publisher
%%  default     : one text column, 10 point font, single spaced (default).
%%  manuscript  : one text column, 12 point font, double spaced article.
%%  preprint    : one text column, 12 point font, single spaced article.  
%%  preprint2   : two text columns, 12 point font, single spaced article.
%%  modern      : a stylish, single text column, 12 point font, article with
%% 		  wider left and right margins. This uses the Daniel
%% 		  Foreman-Mackey and David Hogg design.
%%
%% Note that you can submit to the AAS Journals in any of these 6 styles.
%%
%% There are other optional arguments one can invoke to allow other stylistic
%% actions. The available options are:
%%
%%   astrosymb    : Loads Astrosymb font and define \astrocommands. 
%%   tighten      : Makes baselineskip slightly smaller, only works with 
%%                  the twocolumn substyle.
%%   times        : uses times font instead of the default.
%%   linenumbers  : turn on linenumbering. Note this is mandatory for AAS
%%                  Journal submissions and revisions.
%%   trackchanges : Shows added text in bold.
%%   longauthor   : Do not use the more compressed footnote style (default) for 
%%                  the author/collaboration/affiliations. Instead print all
%%                  affiliation information after each name. Creates a much 
%%                  longer author list but may be desirable for short 
%%                  author papers.
%% twocolappendix : make 2 column appendix.
%%   anonymous    : Do not show the authors, affiliations, acknowledgments,
%%                  and author contributions for dual anonymous review.
%%  resetfootnote : Reset footnotes to 1 in the body of the manuscript.
%%                  Useful when there are a lot of authors and affiliations
%%		    in the front matter.
%%   longbib      : Print article titles in the references. This option
%% 		    is mandatory for PSJ manuscripts.
%%
%% Since v6, AASTeX has included \hyperref support. While we have built in 
%% specific %% defaults into the classfile you can manually override them 
%% with the \hypersetup command. For example,
%%
%% \hypersetup{linkcolor=red,citecolor=green,filecolor=cyan,urlcolor=magenta}
%%
%% will change the color of the internal links to red, the links to the
%% bibliography to green, the file links to cyan, and the external links to
%% magenta. Additional information on \hyperref options can be found here:
%% https://www.tug.org/applications/hyperref/manual.html#x1-40003
%%
%% The "bookmarks" has been changed to "true" in hyperref
%% to improve the accessibility of the compiled pdf file.
%%
%% If you want to create your own macros, you can do so
%% using \newcommand. Your macros should appear before
%% the \begin{document} command.
%%
\newcommand{\vdag}{(v)^\dagger}
\newcommand\aastex{AAS\TeX}
\newcommand\latex{La\TeX}

\newcommand{\tocite}{\textcolor{red}{CITE} }
\newcommand{\TAO}[1]{\textcolor{purple}{[tao:#1]}}
% \newcommand{\edit}[1]{\textbf{#1}}
\newcommand{\edit}[1]{#1}

\newcommand{\ICO}{{I_{\rm CO}}}
\newcommand{\Ijwst}{{I_{\rm JWST}}}
\newcommand{\Icont}{{I_{\rm F200W}}}
\newcommand{\Iblueraw}{{I_{\rm F770W}}}
\newcommand{\Iblue}{{I_{\rm F770W, PAH}}}
\newcommand{\Igreen}{{I_{\rm F1130W}}}
\newcommand{\Ired}{{I_{\rm F2100W}}}

\newcommand{\blueband}{{\rm F770W}_{\rm PAH}}
\newcommand{\greenband}{{\rm F1130W}}
\newcommand{\redband}{{\rm F2100W}}

\newcommand{\betablue}{\beta_{\blueband}}
\newcommand{\betagreen}{\beta_{\greenband}}
\newcommand{\betared}{\beta_{\redband}}
\newcommand{\betaCO}{\beta_{CO}}


\newcommand{\nPAH}{n_{\rm PAH}}
\newcommand{\ndust}{n_{\rm dust}}
\newcommand{\ngas}{n_{\rm gas, cold}}

\newcommand{\kKS}{k_{\rm KS}}
\newcommand{\bKS}{b_{\rm KS}}
\newcommand{\sigKS}{\sigma_{\rm KS}}

\newcommand{\Hmolecular}{H_2}
\newcommand{\HII}{H~{\sc ii}}
\newcommand{\HI}{H~{\sc i}}
\newcommand{\Halpha}{H$\alpha$~$\lambda$6563}
\newcommand{\Hbeta}{H$\beta$~$\lambda$4861}
\newcommand{\OIII}{[O~{\sc iii}]~$\lambda$5007}
\newcommand{\SII}{[S~{\sc ii}]~$\lambda$$\lambda$6717,6731}
\newcommand{\NII}{[N~{\sc ii}]~$\lambda$6583}

\newcommand{\Mstar}{M_{\star}}
\newcommand{\SFR}{{\rm SFR}}
\newcommand{\sSFR}{{\rm sSFR}}
\newcommand{\LCO}{L_{\rm CO}}
\newcommand{\MHI}{M_{\rm H I}}
\newcommand{\FWHMa}{{\rm FWHM}_{\rm ang.}}
\newcommand{\FWHMp}{{\rm FWHM}_{\rm phy.}}
\newcommand{\lowb}{low-$b$}
\newcommand{\highb}{high-$b$}


\newcommand{\btheta}{\boldsymbol {\theta}}
\newcommand{\bphi}{\boldsymbol {\phi}}
\newcommand{\bx}{\mathbf{x}}
\newcommand{\sigmain}{\sigma_{in}}
\newcommand{\err}[2][]{#2_{\text{err}#1}}
\newcommand{\bxerr}{\err{\bx}}
\newcommand{\bxerri}{\err[,i]{\bx}}
\newcommand{\yerr}{\err{y}}
\newcommand{\yerri}{\err[,i]{y}}
\newcommand{\obs}[2][]{#2_{\text{obs}#1}}
\newcommand{\bxobs}{\obs{\bx}}
\newcommand{\bxobsi}{\obs[,i]{\bx}}
\newcommand{\yobsi}{\obs[,i]{y}}


\newcommand{\reproject}{{\tt reproject}}
\newcommand{\linmix}{{\tt LINMIX}}
\newcommand{\mlinmix}{{\tt mLINMIX}}
\newcommand{\modr}{{\tt mODR}}
\newcommand{\mlmethod}{{ML based method}}
\newcommand{\ksmethod}{{KS-test based method}}
\newcommand{\raddest}{{\tt raddest}}

\newcommand{\python}{{\tt Python}}
\newcommand{\jax}{{\tt JAX}}

\newcommand{\tdks}{2D KS test}
\newcommand{\ctest}{Cram\'er test}

\usepackage{amsmath}
\usepackage{multirow}
\usepackage{enumitem}

\defcitealias{Jing2024}{JL25}

%%%%%%%%%%%%%%%%%%%%%%%%%%%%%%%%%%%%%%%%%%%%%%%%%%%%%%%%%%%%%%%%%%%%%%%%%%%%%%%%
%%
%% The following section outlines numerous optional output that
%% can be displayed in the front matter or as running meta-data.
%%
%% Running header information. A short title on odd pages and 
%% short author list on even pages. Note that this
%% information may be modified in production.
\shorttitle{CO-MIR Correlations}
\shortauthors{Jing \& Li}
%%
%% Include dates for submitted, revised, and accepted.
%%\received{February 1, 2025}
%%\revised{March 1, 2025}
%%\accepted{\today}
%%
%% Indicate AAS Journal the manuscript was submitted to.
%%\submitjournal{PSJ}
%% Note that this command adds "Submitted to " the argument.
%%
%% You can add a light gray and diagonal water-mark to the first page 
%% with this command:
%% \watermark{text}
%% where "text", e.g. DRAFT, is the text to appear.  If the text is 
%% long you can control the water-mark size with:
%% \setwatermarkfontsize{dimension}
%% where dimension is any recognized LaTeX dimension, e.g. pt, in, etc.
%%%%%%%%%%%%%%%%%%%%%%%%%%%%%%%%%%%%%%%%%%%%%%%%%%%%%%%%%%%%%%%%%%%%%%%%%%%%%%%%
%%
%% Use this command to indicate a subdirectory where figures are located.
\graphicspath{{./}{figures/}}
%% This is the end of the preamble.  Indicate the beginning of the
%% manuscript itself with \begin{document}.

\begin{document}

\title{Correlations of ALMA CO(2-1) with JWST mid-infrared fluxes down to scale of $\lesssim$100 parsec in nearby star-forming galaxies from PHANGS}

%% A significant change from AASTeX v6+ is in the author blocks. Now an email
%% address is required for each author. This means that each author requires
%% at least one of the following:
%%
%% \author
%% \affiliation
%% \email
%%
%% If these three commands are not available for each author, the latex
%% compiler will issue an error and if you force the latex compiler to continue,
%% it will generate an incomplete pdf.
%%
%% Multiple \affiliation commands are allowed and authors can also include
%% an optional \altaffiliation to indicate a status, i.e. Hubble Fellow. 
%% while affiliations are indexed as footnotes, altaffiliations are noted with
%% with a non-numeric footnote that is set away from the numeric \affiliation 
%% footnotes. NOTE that if an \altaffiliation command is used it must 
%% come BEFORE the \affiliation call, right after the \author command, in 
%% order to place the footnotes in the proper location. Because non-numeric
%% symbols are used, \altaffiliation should be used sparingly.
%%
%% In v7 the \author command takes an optional argument which provides 
%% additional metadata about the author. Authors can provide the 16 digit 
%% ORCID, the surname (family or last) name, the given (first or fore-) name, 
%% and a name suffix, e.g. "Jr.". The syntax is:
%%
%% \author[orcid=0000-0002-9072-1121,gname=Gregory,sname=Schwarz]{Greg Schwarz}
%%
%% This name metadata in not shown, it is only for parsing by the peer review
%% system so authors can be more easily identified. This name information will
%% also be sent to the publisher so they can include it in the CROSSREF 
%% metadata. Including an orcid will hyperlink the author name to the 
%% author's ORCID page. Note that  during compilation, LaTeX will do some 
%% limited checking of the format of the ID to make sure it is valid. If 
%% the "orcid-ID.png" image file is  present or in the LaTeX pathway, the 
%% ORCID icon will appear next to the authors name.
%%
%% Even though emails are now required for each author, the \email does not
%% produce output in the compiled manuscript unless the optional "show" command
%% is used. For example,
%%
%% \email[show]{greg.schwarz@aas.org}
%%
%% All "shown" emails are show in the bottom left of the first page. Due to
%% space constraints, only a few emails should be shown. 
%%
%% To identify a corresponding author, use the \correspondingauthor command.
%% The command appends "Corresponding Author: " to the argument it appears at
%% the bottom left of the first page like the output from \email. 


\author[0009-0004-6271-4321]{Tao Jing}
\affiliation{Department of Astronomy, Tsinghua University, Beijing 100084, China}
\email[show]{tao.jing.astro@gmail.com}

\author[0000-0002-8711-8970]{Cheng Li}
\affiliation{Department of Astronomy, Tsinghua University, Beijing 100084, China}
\email[show]{cli2015@tsinghua.edu.cn}


\correspondingauthor{Tao Jing, Cheng Li}


%% Use the \collaboration command to identify collaborations. This command
%% takes an optional argument that is either a number or the word "all"
%% which tells the compiler how many of the authors above the command to
%% show. For example "\collaboration[all]{(DELVE Collaboration)}" wil include
%% all the authors above this command.
%%
%% Mark off the abstract in the ``abstract'' environment. 
\begin{abstract}
We investigate the correlations of CO (2-1) emission ($\ICO$) with PAH ($\Iblue$ and $\Igreen$) and dust ($\Ired$) emission down to scales of $\lesssim$ 100 pc, by applying \raddest, a novel regression technique recently developed by \citet{Jing2024} that effectively handles uncertainties and outliers in datasets, to 19 nearby star-forming galaxies in the PHANGS sample. We find that for the majority of the data points in all galaxies, the scaling of $\ICO$ with $\Iblue$, $\Igreen$, and $\Ired$ can be well described by log-log linear relations, though with substantial dependence on ionization conditions (i.e., \HII-like, composite-like, and AGN-like). Under given ionization conditions, significant galaxy-to-galaxy variations are identified, and are primarily attributed to variations of intercept $b$, which exhibits clear bimodality. \edit{This bimodality is related to the normalized overall host galaxy star formation rate, such as specific star formation and star formation efficiency.} The differences in slope $k$ and intrinsic scatter $\sigma$ across different MIR bands ($\Iblue$, $\Igreen$, and $\Ired$) are minor compared to their galaxy-to-galaxy variations. All parameters ($k$, $b$, and $\sigma$) depend on the spatial scale of measurement, suggesting that the coupling among CO, PAH, and dust is regulated by different mechanisms at varying scales. \edit{We identify deviations from the log–log linear relation in the brightest regions, primarily characterized by a flattening of the slope.} No significant (3$\sigma$) correlations are found between global properties and the best-fit parameters. We discuss the comparison to previous studies and plausible physics behind the statistical results obtained in this work. %We suggest that the obtained results should be compared with simulations incorporating detailed molecular gas, dust, and PAH models to release their full potential on constraining the corresponding physical processes.
\end{abstract}

%% Keywords should appear after the \end{abstract} command. 
%% The AAS Journals now uses Unified Astronomy Thesaurus (UAT) concepts:
%% https://astrothesaurus.org
%% You will be asked to selected these concepts during the submission process
%% but this old "keyword" functionality is maintained in case authors want
%% to include these concepts in their preprints.
%%
%% You can use the \uat command to link your UAT concepts back its source.
\keywords{}

%% From the front matter, we move on to the body of the paper.
%% Sections are demarcated by \section and \subsection, respectively.
%% Observe the use of the LaTeX \label
%% command after the \subsection to give a symbolic KEY to the
%% subsection for cross-referencing in a \ref command.
%% You can use LaTeX's \ref and \label commands to keep track of
%% cross-references to sections, equations, tables, and figures.
%% That way, if you change the order of any elements, LaTeX will
%% automatically renumber them.

\section{Introduction} \label{sec:introduction}

...


\end{document}

% End of file `sample7.tex'.
